% Solutions to the exercises of Chapter 3.
\begin{solutions}
  \item \emph{Without pivoting} the multiplier is $\ell_{21} = 1/10^{-20} = 10^{20}$. The second row becomes
        \[
          (1 - 10^{20})\,x_2 = 2 - 10^{20},
        \]
        and in double precision both $1 - 10^{20}$ and $2 - 10^{20}$ round to $-10^{20}$: the information carried by the original entries
        $1$ and $2$ is lost. Hence $x_2 = 1$ and back substitution gives $x_1 = (1 - x_2)/10^{-20} = 0$. \emph{With partial pivoting} the rows
        are swapped, the multiplier is $10^{-20}$, the second row becomes $(1 - 10^{-20})\,x_2 = 1 - 2\times10^{-20}$, which rounds to
        $x_2 = 1$, and $x_1 = 2 - x_2 = 1$. The exact solution is $x_1 = 1/(1 - 10^{-20})$, $x_2 = (1 - 2\times10^{-20})/(1 - 10^{-20})$, both equal to 1
        to all printed digits. The matrix is well conditioned ($\kappa \approx 2.6$); the failure is entirely the algorithm's. A tiny pivot
        produces a huge multiplier, the elimination step adds huge numbers to small ones, and the rounding errors of that step are
        of the size of the huge numbers. Pivoting keeps every multiplier at most one in magnitude, which bounds this growth.

  \item Column by column with the formulas of \cref{thm:cholesky}: $L_{11} = \sqrt4 = 2$, $L_{21} = 2/2 = 1$, $L_{31} = 2/2 = 1$;
        $L_{22} = \sqrt{5 - 1^2} = 2$, $L_{32} = (3 - 1\cdot1)/2 = 1$; $L_{33} = \sqrt{6 - 1^2 - 1^2} = 2$. So
        \[
          L = \begin{bmatrix}2 & 0 & 0\\ 1 & 2 & 0\\ 1 & 1 & 2\end{bmatrix}.
        \]
        Forward substitution $Ly = (8, 10, 11)$ gives $y = (4, 3, 2)$, and back substitution $L\T x = y$ gives $x_3 = 1$, $x_2 = (3 - 1)/2 = 1$,
        $x_1 = (4 - 1 - 1)/2 = 1$. Check: $S(1, 1, 1)\T = (8, 10, 11)\T$.

  \item The block matrix is a sum of PSD terms:
        \[
          M = \begin{bmatrix}S & HP\\ PH\T & P\end{bmatrix}
            = \begin{bmatrix}H\\ I\end{bmatrix}P\begin{bmatrix}H\T & I\end{bmatrix} + \begin{bmatrix}R & 0\\ 0 & 0\end{bmatrix} \succeq 0,
        \]
        the first a congruence of $P \succ 0$ and the second PSD because $R \succ 0$ (\cref{prop:psd}). Since $S = HPH\T + R \succ 0$, the block
        factorization of \cref{prop:block-inverse} writes $M$ as a congruence of $\diag(S,\ M/S)$ by an invertible unit-triangular matrix, with
        $M/S = P - PH\T S\inv HP$. A congruence by an invertible matrix preserves semi-definiteness in both directions, so $M \succeq 0$ implies
        $M/S \succeq 0$. Finally $KSK\T = PH\T S\inv SS\inv HP = PH\T S\inv HP$, so $P - KSK\T = M/S \succeq 0$. This is the statement that a Kalman
        update never increases uncertainty (\cref{ch:kf}).

  \item A compact implementation, storing each reflection and applying it to the remaining columns:
\begin{lstlisting}[style=py]
def householder_qr(A):
    A = A.astype(float).copy(); m, n = A.shape; vs = []
    for k in range(n):
        x = A[k:, k].copy()
        v = x.copy(); v[0] += np.copysign(np.linalg.norm(x), x[0])
        v /= np.linalg.norm(v)
        A[k:, k:] -= 2.0 * np.outer(v, v @ A[k:, k:])   # apply H_k
        vs.append(v)
    Q = np.eye(m)
    for k in reversed(range(n)):                        # Q = H_1 H_2 ... H_n
        Q[k:, :] -= 2.0 * np.outer(vs[k], vs[k] @ Q[k:, :])
    return Q, np.triu(A[:n, :])
\end{lstlisting}
        For a random $10\times3$ matrix, $\|Q\T Q - I\| \approx 10^{-15}$, $\|Q_1R_1 - A\| \approx 2\times10^{-15}$ and $\|Q_2\T A\| \approx 5\times10^{-16}$: the last
        seven columns of $Q$ are an orthonormal basis of the left nullspace, as \cref{prop:qr-nullspace} states. For the least-squares test, build
        $A = U\diag(1, 10^{-4}, 10^{-8})V\T$ with random orthonormal $U \in \R^{10\times3}$ and $V$, choose $x = (1, 2, 3)$ and $b = Ax$. A typical run
        gives a relative error of $0.14$ for the normal equations and $3\times10^{-9}$ for QR. The normal equations work with
        $\kappa(A\T A) = 10^{16}$ and lose all sixteen digits; QR works with $\kappa(A) = 10^8$ and keeps about eight, the most any algorithm can
        keep for this problem.

  \item Let $H_f = Q\begin{bmatrix}R_1\\0\end{bmatrix} = Q_1R_1$ with $Q = [Q_1\ Q_2]$ orthogonal and $R_1$ invertible (full column rank). Then
        \[
          H_f(H_f\T H_f)\inv H_f\T = Q_1R_1\big(R_1\T Q_1\T Q_1R_1\big)\inv R_1\T Q_1\T = Q_1R_1R_1\inv R_1\mT R_1\T Q_1\T = Q_1Q_1\T .
        \]
        Because $Q$ is orthogonal, $QQ\T = Q_1Q_1\T + Q_2Q_2\T = I$, so $I - H_f(H_f\T H_f)\inv H_f\T = Q_2Q_2\T$. Any other orthonormal basis $A_\perp$ of
        $\nullsp(H_f\T) = \range(Q_2)$ (\cref{prop:qr-nullspace}) is $A_\perp = Q_2O$ for an orthogonal matrix $O$, hence
        $A_\perp A_\perp\T = Q_2OO\T Q_2\T = Q_2Q_2\T$, which proves \eqref{eq:nullspace-projector}.

  \item Apply \cref{thm:woodbury} with $A = P\inv$, $U = H\T$, $C = R\inv$, $V = H$:
        \[
          \big(P\inv + H\T R\inv H\big)\inv = P - PH\T\big(R + HPH\T\big)\inv HP .
        \]
        With $n = 75$ states and $m$ measurements, the right side forms $PH\T$ and $HPH\T$ (about $2mn^2 + 2m^2n$ flops), factors the $m\times m$
        matrix $S$ ($m^3/3$) and forms the product (about $2mn^2$); the left side needs $P\inv$, the product $H\T R\inv H$ (about $mn^2$ for a
        diagonal $R$) and the inverse of an $n\times n$ matrix (two $n\times n$ inversions, about $2n^3 \approx 8\times10^5$ flops). For $m = 2$ the
        right side costs about $5\times10^4$ flops, some fifteen times less: this is why the Kalman filter is written in covariance form. For
        $m = 500$ the right side costs about $8\times10^7$ flops (dominated by $m^3/3$ and $2m^2n$), while the left side costs only about
        $4\times10^6$: with more measurements than states, the information form is cheaper. The MSCKF obtains the same saving without leaving
        the covariance form by compressing tall measurements with QR (\cref{prop:compression}).

  \item Let $J = \begin{bmatrix}0 & -1\\ 1 & 0\end{bmatrix}$, so the matrix is $\theta J$. Since $J^2 = -I$, the powers repeat with period four:
        $J^{2k} = (-1)^kI$ and $J^{2k+1} = (-1)^kJ$. Splitting the series into even and odd terms,
        \[
          \exp(\theta J) = \sum_k\frac{(-1)^k\theta^{2k}}{(2k)!}I + \sum_k\frac{(-1)^k\theta^{2k+1}}{(2k+1)!}J = \cos\theta\,I + \sin\theta\,J
          = \begin{bmatrix}\cos\theta & -\sin\theta\\ \sin\theta & \cos\theta\end{bmatrix},
        \]
        the planar rotation by $\theta$ (\cref{ex:rot2d}): the planar instance of \cref{prop:exp-skew}.

  \item $A$ and $B$ are nilpotent ($A^2 = B^2 = 0$), so $\exp(A) = I + A = \begin{bmatrix}1 & 1\\0 & 1\end{bmatrix}$ and
        $\exp(B) = \begin{bmatrix}1 & 0\\1 & 1\end{bmatrix}$, whose product is $\begin{bmatrix}2 & 1\\ 1 & 1\end{bmatrix}$. On the other hand
        $C = A + B = \begin{bmatrix}0 & 1\\1 & 0\end{bmatrix}$ satisfies $C^2 = I$, so
        \[
          \exp(A + B) = \cosh(1)\,I + \sinh(1)\,C = \begin{bmatrix}1.543 & 1.175\\ 1.175 & 1.543\end{bmatrix} \ne \begin{bmatrix}2 & 1\\ 1 & 1\end{bmatrix}.
        \]
        The matrices do not commute: $AB = \diag(1, 0)$ while $BA = \diag(0, 1)$, so property 3 of \cref{prop:matexp} does not apply.

  \item Doubles in $[2^{30}, 2^{31})$ are spaced by $2^{30-52} = 2^{-22}$~s $\approx 238$~ns. Unix time today ($\approx 1.8\times10^9$~s) lies in that
        range. A single $\Dt$ computed as a difference of two such numbers is wrong by up to about $10^{-7}$~s, a relative error of a few
        times $10^{-5}$ on a 5~ms step; at 1~rad/s that is an attitude error of $10^{-7}$~rad per step, well below the gyro noise per step
        ($\sigma_g\sqrt\Dt \approx 10^{-5}$~rad). So differences would be tolerable, but the representation still fails the project's needs:
        \begin{itemize}[leftmargin=1.2em]
          \item it cannot hold EuRoC's nanosecond timestamps exactly, so equal timestamps may compare unequal and event ordering may change;
          \item \emph{accumulating} time is disastrous: adding $0.005$ to a double near $1.79\times10^9$ rounds the same way at every step, and one
                hour at 200~Hz ($720\,000$ additions) drifts by 82~ms in a simulation.
        \end{itemize}
        Absolute time therefore stays in 64-bit integer nanoseconds, and only differences are converted (AGENTS.md C5).

  \item Multiply and divide by the conjugate:
        \[
          \sqrt{x^2 + 1} - x = \frac{1}{\sqrt{x^2 + 1} + x},
        \]
        which adds two positive numbers instead of subtracting nearly equal ones. At $x = 10^8$ the original form returns exactly $0$
        (because $x^2 + 1$ rounds to $x^2$), the rewritten one $5\times10^{-9} \approx 1/(2x)$. Using \code{hypot(x, 1)} for the square root also
        avoids overflow of $x^2$ for $|x| > 10^{154}$. For the second expression use $1 - \cos x = 2\sin^2(x/2)$:
        \[
          \frac{1 - \cos x}{x^2} = \frac12\left(\frac{\sin(x/2)}{x/2}\right)^2,
        \]
        which involves no subtraction (and $\sin(y)/y$ is accurate for small $y$, or can be replaced by its series $1 - y^2/6$). At $x = 10^{-8}$
        the original form returns $0$ and the rewritten one $0.5$; at $x = 10^{-5}$ the original is wrong in the eighth digit
        (\cref{ex:cancellation}).
\end{solutions}
